\begin{thebibliography}{99}
\bibitem{naselaris2011} Naselaris T, Kay KN, Nishimoto S, Gallant JL. Encoding and decoding in fMRI. NeuroImage 2011;56(2):400-410.
\bibitem{yamins2016} Yamins DLK, DiCarlo JJ. Using goal-driven deep learning models to understand sensory cortex. Nature Neuroscience 2016;19(3):356-365.
\bibitem{kragel2019} Kragel PA, Reddan MC, LaBar KS, Wager TD. Emotion schemas are embedded in the human visual system. Science Advances 2019;5(7):eaaw4358.
\bibitem{jang2025} Jang G, Kragel PA. Understanding human amygdala function with artificial neural networks. Journal of Neuroscience 2025;45(18):e1436242025.
\bibitem{kriegeskorte2019} Kriegeskorte N, Douglas PK. Interpreting encoding and decoding models. Current Opinion in Neurobiology 2019;55:167-179.
\bibitem{liu2024} Liu Y. Emergence of emotion selectivity in deep neural networks trained to recognize visual objects. PLoS Computational Biology 2024;20(3):e1011943.
\bibitem{gao2025} Gao S, Ajith A, Peelen MV. Object representations drive emotion schemas across a large and diverse set of daily-life scenes. Communications Biology 2025.
\bibitem{aliko2020} Aliko S, Huang J, Gheorghiu F, Meliss S, Skipper JI. A naturalistic neuroimaging database for understanding the brain using ecological stimuli. Scientific Data 2020;7:347.
\bibitem{toisoul2021} Toisoul A, Kossaifi J, Bulat A, Tzimiropoulos G, Pantic M. Estimation of continuous valence and arousal levels from faces in naturalistic conditions. Nature Machine Intelligence 2021;3:42-50.
\bibitem{wang2017} Wang S, Tudusciuc O, Mamelak AN, Ross IB, Adolphs R, Rutishauser U. The human amygdala parametrically encodes the intensity of specific facial emotions and their categorical ambiguity. Nature Communications 2017;8:14821.
\bibitem{oneill2026} O'Neill PK, Posani L, Salzman CD. Representational geometry of emotional states in basolateral amygdala. Nature Neuroscience 2026.
\bibitem{wagner2023} Wagner J, Triantafyllopoulos A, Wierstorf H, Schmitt M, Burkhardt F, Eyben F, Schuller BW. Dawn of the transformer era in speech emotion recognition. IEEE Transactions on Pattern Analysis and Machine Intelligence 2023;45(9):10745-10759.
\bibitem{sander2003} Sander D, Grafman J, Zalla T. The human amygdala: an evolved system for relevance detection. Reviews in the Neurosciences 2003;14(4):303-316.
\bibitem{dicarlo2026} DiCarlo JJ. An image-computable characterization of the non-conditioned linkage of visual drive and valence in the primate amygdala. bioRxiv 2026.
\bibitem{dupre2022} Dupre la Tour T, Eickenberg M, Nunez-Elizalde AO, Gallant JL. Feature-space selection with banded ridge regression. NeuroImage 2022;264:119728.
\bibitem{nili2014} Nili H, Wingfield C, Walther A, Su L, Marslen-Wilson W, Kriegeskorte N. A toolbox for representational similarity analysis. PLoS Computational Biology 2014;10(4):e1003553.
\bibitem{soderberg2026} Soderberg C, Jang G, Kragel PA. Sensory encoding of emotion conveyed by the face and visual context. Social Cognitive and Affective Neuroscience 2026;21(1):nsag028.
\bibitem{morgenroth2025} Morgenroth E. Emo-FilM: a multimodal dataset for affective neuroscience. Scientific Data 2025;12:684.
\end{thebibliography}
